\renewcommand{\thesection}{\arabic{section}}\renewcommand{\theHsection}{chapter.\arabic{section}}\setcounter{section}{1}\section{守恒定律}
\noindent\begin{minipage}{\textwidth}\ProblemHeading{例题}{2.1（改写）}
一组相互作用的质点处于沿 $z$ 轴方向的匀强外场中，质点间的作用力来自只依赖相对距离的势能。外场作用在各质点上的力均平行于 $z$ 轴，且各自恒定。
\SubQ{1}{根据平移对称性，判断总动量的哪些分量守恒。}
\SubQ{2}{证明对于任意选取的坐标原点，总动量矩在 $z$ 方向的分量均守恒。}
\end{minipage}\par\vfill
\noindent\begin{minipage}{\textwidth}\ProblemHeading{习题}{2.2}
质量为 $m$ 的质点以速度 $\mathbf{v}_1$ 从势能为常数 $U_1$ 的半空间运动到另一个势能为常数 $U_2$ 的半空间。求质点运动方向的改变。
\end{minipage}\par\vfill
\newpage
\noindent\begin{minipage}{\textwidth}\ProblemHeading{习题}{2.3}
两个惯性参考系 $K$ 和 $K\prime$ 的坐标轴彼此平行，$K\prime$ 相对 $K$ 以恒定速度 $\mathbf V$ 运动。求同一个封闭力学系统沿同一段运动在这两个参考系中的作用量之间的关系。
\end{minipage}\par\vfill
\noindent\begin{minipage}{\textwidth}\ProblemHeading{例题}{2.4（改写）}
一组质点在以固定点 $O$ 为中心的外场中运动。各质点的外场势能只依赖其到 $O$ 的距离，质点间的相互作用势能只依赖彼此的距离。利用旋转对称性，证明系统相对于 $O$ 的总动量矩守恒。
\end{minipage}\par\vfill
\newpage
\noindent\begin{minipage}{\textwidth}\ProblemHeading{习题}{2.5}
在柱坐标 $(r,\varphi,z)$ 中，用质点的坐标及其时间导数表示质点相对于原点的动量矩的各笛卡儿分量及其大小，其中 $r$ 为到 $z$ 轴的距离，$\varphi$ 为方位角。
\end{minipage}\par\vfill
\noindent\begin{minipage}{\textwidth}\ProblemHeading{习题}{2.6}
在球坐标 $(r,\theta,\varphi)$ 中，用质点的坐标及其时间导数表示质点相对于原点的动量矩的各笛卡儿分量及其大小，其中 $r$ 为到原点的距离，$\theta$ 为极角，$\varphi$ 为方位角。
\end{minipage}\par\vfill
\newpage
\noindent\begin{minipage}{\textwidth}\ProblemHeading{习题}{2.7}
取适合各场对称性的直角坐标系，说明坐标轴的选取，并判断力学系统在下列场中运动时，动量 $\mathbf P$ 和动量矩 $\mathbf M$ 的哪些分量守恒。
\SubQ{1}{无限大均匀平面场。}
\SubQ{2}{无限长均匀圆柱场。}
\SubQ{3}{无限长均匀棱柱场。}
\end{minipage}\par\vfill
\noindent\begin{minipage}{\textwidth}\ProblemHeading{习题}{2.8}
取适合各场对称性的直角坐标系，说明坐标轴的选取，并判断力学系统在下列场中运动时，动量 $\mathbf P$ 和动量矩 $\mathbf M$ 的哪些分量守恒。
\SubQ{1}{两个点场。}
\SubQ{2}{无限大均匀半平面场。}
\SubQ{3}{均匀圆锥场。}
\end{minipage}\par\vfill
\newpage
\noindent\begin{minipage}{\textwidth}\ProblemHeading{习题}{2.9}
取适合各场对称性的直角坐标系，说明坐标轴的选取，并判断力学系统在下列场中运动时，动量 $\mathbf P$ 和动量矩 $\mathbf M$ 的哪些分量守恒。
\SubQ{1}{均匀圆环场。}
\SubQ{2}{无限长均匀圆柱形螺旋线场。取螺旋轴为 $z$ 轴，沿螺旋线绕该轴正向前进一周时，$z$ 坐标的增量为 $h$。求 $P_z$ 和 $M_z$ 的一个守恒的线性组合。}
\end{minipage}\par\vfill
\noindent\begin{minipage}{\textwidth}\ProblemHeading{例题}{2.10（改写）}
一个质点在稳定平衡位置附近做一维微小振动，势能相对于平衡值的增量是位移的二次齐次函数。用力学相似性判断振动周期是否随振幅改变。
\end{minipage}\par\vfill
\newpage
\noindent\begin{minipage}{\textwidth}\ProblemHeading{例题}{2.11（改写）}
在同一匀强重力场中，两个质点分别从不同高度由静止开始自由下落。用力学相似性求下落时间之比与下落高度之比的关系。
\end{minipage}\par\vfill
\noindent\begin{minipage}{\textwidth}\ProblemHeading{例题}{2.12（改写）}
考虑同一对只受相互牛顿引力作用的质点。若其相对运动可以沿两条几何相似的闭合轨道进行，用力学相似性求两轨道周期之比与轨道线性尺寸之比的关系。
\end{minipage}\par\vfill
\newpage
\noindent\begin{minipage}{\textwidth}\ProblemHeading{习题}{2.13}
质量不同、所处势能函数相同的两个质点沿相同轨迹运动，它们通过对应轨迹段所用的时间满足什么关系？
\end{minipage}\par\vfill
\noindent\begin{minipage}{\textwidth}\ProblemHeading{习题}{2.14（校订）}
将势能乘以一个正的常数后，质点沿相同轨迹运动所用的时间有什么变化？
\end{minipage}\par\vfill
